\begin{question}

  \begin{questionpart}
    On p. 247, Gottfried (1966) states that
    \begin{equation*}
      [x_i, G(\vb{p})] = i\hbar \pdv{G}{p_i},
      [p_i,F(\vb{x})] = -i\hbar \pdv{F}{x_i}
    \end{equation*}
    can be ``easily derived'' from the fundamental commutation
    relations for all functions of $F$ and $G$ that can be expressed
    as power series in their arguments.  Verify this statement.
  \end{questionpart}

  \begin{questionpart}
    Evaluate $[x^2,p^2]$.  Compare your result with the classical
    Poisson bracket $[x^2,p^2]_{\text{classical}}$.
  \end{questionpart}
\end{question}

\begin{purpose}
  Formalize a bit further what we already did in problem
  \ref{1.30:theQuestion}.  This one might actually be easy... mostly
  because we've mostly already done it.
\end{purpose}

\begin{answer}

  \begin{answerpart}{``Easily''}

    See the derivation of equation \ref{1.30:xfp} for
    \begin{equation}
      [x_i, G(\vb{p})] = -i\hbar \pdv{G}{p_i}
    \end{equation}


    For the converse, we draw upon our previous finding in equation
    \ref{1.30:IHaveThePower}

    \begin{equation}
      [A,B^n]
      =
      B^{n-1}[A,B]
      + \sum_{i=1}^{n-1}B^{i-1}[A,B]B^{n-i}
    \end{equation}

    Whenever we have a fundamental commutation relationship to draw
    upon which is a constant, we can refer back to equation
    \ref{1.30:ConstantPower} and refine it further

    \begin{equation}
      \inlabel{ABTesting}
      \begin{split}
        [A,B^n]
        &=
        B^{n-1}[A,B]
        + \sum_{i=1}^{n-1}B^{i-1}[A,B]B^{n-i}
        \\
        &=
        [A,B]nB^{n-1}
        \\
      \end{split}
    \end{equation}

    Putting that another way

    \begin{equation}
      [p, x^n] = -i\hbar n x^{n-1} = -i\hbar \pdv{}{x} x^n
    \end{equation}

    Substituting that back into our taylor series just like we did in
    equations \ref{1.30:presub} and \ref{1.30:postsub} and we,
    entirely expectedly, get the same result (just with a negative
    sign from reversing the order of the canonical commutation and a
    different variable we're differentiating against).

    So...Is it easily derived?  Yeah...if you've already derived it.
  \end{answerpart}

  \begin{answerpart}{``Derived''}

    It is tempting to try pulling in equation \inref{ABTesting}

    \begin{equation}
      \begin{alignedat}{2}
        && [x^2,p^2]
        &= [x^2,p]\pdv{}{p}p^2
        \\
        &&&= -[p,x^2]\pdv{}{p}p^2\\
        &&&= -[p,x]\pdv{}{}x^2\pdv{}{p}p^2\\
        &&&= -i\hbar\pdv{}{x}x^2\pdv{}{p}p^2\\
        &&&= -i4\hbar x p\\
      \end{alignedat}
    \end{equation}

    Buuuuut we can't do that because equation \inref{ABTesting}
    assumes that $[A,B] = k$ whereas $[x^2, p]$ does \emph{not}
    produce a constant.  Let's try our canonical commutations instead.
    Pulling in the right one from the book

    \begin{equation}
      \tag{1.232e}
      [A,BC] = [A,B]C + B[A,C]
    \end{equation}

    \begin{equation}
      \begin{split}
        [x^2,p^2]
        &= [x^2, pp]\\
        &= [x^2,p]p + p[x^2,p]\\
        &= -([p,x^2]p + p[p,x^2])\\
      \end{split}
    \end{equation}

    Solving for that commutator

    \begin{equation}
      \begin{split}
        [p,x^2] &= [p,x]x + x[p,x]\\
      \end{split}
    \end{equation}

    and plugging it back in
    \begin{equation}
      \begin{split}
        [x^2,p^2]
        &= -(([p,x]x + x[p,x])p + p([p,x]x + x[p,x]))\\
        &= ([x,p]x + x[x,p])p + p([x,p]x + x[x,p])\\
        &= [x,p]xp + x[x,p]p + p[x,p]x + px[x,p]\\
        &= [x,p](xp + xp + px + px)\\
        &= i\hbar(2xp + 2px)\\
        &= i\hbar2\{x,p\}\qed\\
      \end{split}
    \end{equation}

    Pulling in the definition of Poisson brackets from the book
    
    \begin{equation}
      \tag{1.230}
      [A(q,p),B(q,p)]_{\text{classical}} \equiv
      \sum_s \left(
      \pdv{A}{q_s}
      \pdv{B}{p_s}
      - 
      \pdv{A}{p_s}
      \pdv{B}{q_s}
      \right)
    \end{equation}

    and applying it to our scenario
    
    \begin{equation}
      \begin{split}
        [x^2,p^2]_{\text{classical}}
        &\equiv
        \sum_s \left(
        \pdv{x^2}{x}
        \pdv{p^2}{p}
        - 
        \pdv{x^2}{p}
        \pdv{p^2}{x}
        \right)
        \\
        &=
        \pdv{x_i^2}{x_i}
        \pdv{p_i^2}{p_i}
        -
        \underbrace{
          \pdv{x_i^2}{p_i}
          \pdv{p_i^2}{x_i}
        }_{0}
        \\
        &= 4x_ip_i
        \\
        &= (2x_ip_i + 2x_ip_i)
        \\
        &= \frac{1}{i\hbar}[x^2,p^2]_{\text{quantum}}\qed
        \\
      \end{split}
    \end{equation}

  \end{answerpart}

\end{answer}
